% Generated draft: requires mathematical and source review.
\documentclass[11pt]{article}
\usepackage[margin=25mm]{geometry}
\usepackage{fontspec}
\setmainfont{Latin Modern Roman}
\usepackage{amsmath,amssymb,mathtools,unicode-math}
\setmathfont{Latin Modern Math}
\usepackage{enumitem,needspace,xurl,hyperref}
\setlength{\parindent}{0pt}
\setlength{\parskip}{5pt}
\setlength{\emergencystretch}{3em}
\newcommand{\N}{\mathbb{N}}
\newcommand{\Z}{\mathbb{Z}}
\newcommand{\Q}{\mathbb{Q}}
\newcommand{\R}{\mathbb{R}}
\newcommand{\C}{\mathbb{C}}
\newcommand{\F}{\mathbb{F}}
\newcommand{\A}{\mathbb{A}}
\newcommand{\PP}{\mathbb{P}}
\newcommand{\E}{\mathbb{E}}
\newcommand{\eps}{\varepsilon}
\newcommand{\dd}{\,\mathrm{d}}
\begin{document}
\section*{UnsolvedMath ID: 30005358}
\textbf{Source problem code:} catalogue entry
\subsection{Definitions and mathematical statement}
The supplied question is a research programme rather than a specified conjecture. The following gives a standard framework in which to state it; the source does not uniquely specify the intended flow or invariant.

Let \(M\preccurlyeq N\) be an elementary extension of first-order structures, and let \(G\) be an infinite group definable over \(M\). Put \(G_0=G(M)\) and \(H=G(N)\). Let \(\operatorname{Def}_M(G_0)\) be the Boolean algebra of subsets of \(G_0\) definable with parameters from \(M\). For a Boolean algebra \(\mathcal A\) of subsets of a group \(K\), write \(S(\mathcal A)\) for its Stone space of ultrafilters, with basic clopen sets \([U]=\{p:U\in p\}\). If \(\mathcal A\) is invariant under left translations, the action is defined by
\[
 U\in g\cdot p\quad\Longleftrightarrow\quad g^{-1}U\in p.
\]
In the standard convention, \(\mathcal A\) is \(d\)-closed if, for every \(p\in S(\mathcal A)\) and \(U\in\mathcal A\),
\[
 d_pU:=\{g\in K:g^{-1}U\in p\}\in\mathcal A.
\]
The supplied excerpt does not define its \(d\)-operator; agreement with this convention requires review.

For a compact \(K\)-flow \(X\), its Ellis semigroup is
\[
 E_K(X)=\overline{\{x\mapsto g\cdot x:g\in K\}}\subseteq X^X,
\]
with the product topology and composition of maps. If \(I\) is a minimal left ideal and \(u\in I\) is idempotent, \(uI\) is an Ellis group. These definitions use ordinary composition; the excerpt instead describes a left-continuous multiplication on \(S(\mathcal A)\), whose convention is not recoverable from its corrupted product symbol.

\textbf{Research question.} Determine how the definable topological dynamics associated with \(G(M)\) changes on passage to \(G(N)\), for elementary extensions \(M\preccurlyeq N\). In particular, describe the relationship between the corresponding flows and Ellis groups. The source does not specify which flows, which comparison maps, which group topology, or which additional hypotheses are intended, so it does not determine a single theorem to prove.

\textbf{Supplied contextual result.} The adjacent excerpt attributed to Ludomir Newelski, in joint work with Adam Malinowski, considers infinite group structures \(G_0\) and \(H\), a \(d\)-closed translation-invariant Boolean subalgebra \(\mathcal A\subseteq\operatorname{Def}(G_0)\), and a \(d\)-closed translation-invariant Boolean algebra \(\mathcal B\) on \(H\), subject to
\[
 U(H)\in\mathcal B\quad(U\in\mathcal A),\qquad
 \mathcal B\mathbin{\upharpoonright}G_0
 :=\{V\cap G_0:V\in\mathcal B\}=\mathcal A.
\]
Here \(U(H)\) denotes interpretation in the extension of a formula defining \(U\); this notation presupposes a compatible elementary-extension setting. The relation between the group structures is corrupted in the original excerpt, so this reconstruction is qualified by the clean question's explicit reference to elementary extension.

The excerpt asserts that \(S(\mathcal A)\), with its unspecified multiplication, is a left-continuous semigroup isomorphic to its Ellis semigroup. It further reports: if there are generic points in \(S(\mathcal B)\), then every Ellis group \(K_{\mathcal A}\) of \(S(\mathcal A)\) is a homomorphic image of a subgroup of some Ellis group \(K_{\mathcal B}\) of \(S(\mathcal B)\); symbolically,
\[
 \forall K_{\mathcal A}\ \exists K_{\mathcal B}\ \exists L\leq K_{\mathcal B}\ \exists\text{ a surjective group homomorphism }L\longrightarrow K_{\mathcal A}.
\]
The excerpt does not define generic points or state continuity properties of this homomorphism. This is a supplied conditional result, not a conjecture asserted here or a resolution of the broader question.
\subsection{Short English statement}
Describe how the definable flows and Ellis groups associated with an infinite definable group change under elementary extension. The source gives a broad question and a conditional comparison result involving generic points, but leaves the precise comparison target and several conventions unresolved.
\subsection{Sources}
\begin{itemize}
\item Local registry: \texttt{data/unsolvedmath/problems.json}, permanent UnsolvedMath ID \(30005358\), source problem code \texttt{catalogue entry}, title \emph{Elementary Extensions of Definable Group Dynamics}. The supplied record labels the problem open and its statement \texttt{corrected\_verified}; these are supplied curation labels, not independent verification.
\item \emph{Model Theory: Combinatorics, Groups, Valued Fields and Neostability} (2023), Oberwolfach Reports, pp.\ 91--149, DOI 10.4171/owr/2023/2: \url{https://doi.org/10.4171/owr/2023/2}. This is also the supplied literature link labelled ``original Oberwolfach report.'' The adjacent ``Ellis groups'' excerpt names Ludomir Newelski, with joint work by Adam Malinowski; the record does not identify a proposer of the broad question.
\item Supplied literature reference: \emph{Topological dynamics and NIP fields}, \url{https://doi.org/10.1016/j.apal.2021.103010}.
\item Supplied literature reference labelled \emph{Ellis group conjecture and variants}, \url{https://www.cambridge.org/core/services/aop-cambridge-core/content/view/D86369293DF898D468876576AD24F200/S0022481218000555a.pdf}.
\item Literature assessment dated exactly \texttt{2026-08-22}: the supplied curation labels the question open and says that later work supplies model-independence or homomorphism results under extra hypotheses, as well as counterexamples to the general Ellis-group conjecture, without establishing a general resolution. This is supplied dated curation derived from the OWR material, not an independently verified literature assessment; its status and cited claims require independent checking.
\end{itemize}
\end{document}
